\documentclass[10pt]{article}
\usepackage[T1]{fontenc}
\usepackage[utf8]{inputenc}
\usepackage[french]{babel}
\usepackage{amsmath,amssymb,mathrsfs,mathpazo}
\usepackage{geometry}
\geometry{a4paper,margin=2cm}
\usepackage{graphicx,array,tabularx,tabularray,colortbl,multirow,multicol,booktabs,enumitem,pifont,ulem,fancybox,diagbox,cancel}
\usepackage{pstricks,pstricks-add,pst-plot,pst-tree,pst-node,pst-text,pst-func,pst-math,pst-bspline,pst-3dplot,pst-eucl,pst-solides3d}
\let\pscalculate\undefined
\usepackage{pst-bezier}
\usepackage{tikz,pgfplots,tkz-tab}
\usetikzlibrary{arrows,arrows.meta,patterns,calc,positioning,trees,decorations.markings,shapes.geometric,shapes.misc,shapes.arrows,intersections,decorations.pathmorphing,decorations.pathreplacing,backgrounds,shadows,babel}
\pgfplotsset{compat=1.18}
\usepgfplotslibrary{fillbetween}
\usepackage[np]{numprint}
\usepackage{listings,makecell,pas-tableur,twoopt}
\let\barre\undefined
\usepackage[pstricks]{tablvar}
\definecolor{darkgreen}{RGB}{0,100,0}
\providecommand{\red}{\color{red}}
\providecommand{\blue}{\color{blue}}
\providecommand{\green}{\color{green}}
\providecommand{\black}{\color{black}}
\providecommand{\white}{\color{white}}
\providecommand{\gray}{\color{gray}}
\providecommand{\cyan}{\color{cyan}}
\providecommand{\magenta}{\color{magenta}}
\providecommand{\yellow}{\color{yellow}}
\providecommand{\violet}{\color{violet}}
\providecommand{\orange}{\color{orange}}
\providecommand{\euro}{\texteuro}
\providecommand{\EUR}{\texteuro}
\providecommand{\ieme}{\textsuperscript{e}}
\providecommand{\degres}{\ensuremath{^\circ}}
\providecommand{\degree}{\ensuremath{^\circ}}
\providecommand{\vv}[1]{\overrightarrow{#1}}
\providecommand{\Coord}[1]{#1}
\providecommand{\decosix}{}
\providecommand{\starredbullet}{\ensuremath{\bullet}}
\pagestyle{empty}
\setlength{\parindent}{0pt}
\usepackage{adjustbox}\DeclareUnicodeCharacter{2236}{:}\DeclareUnicodeCharacter{2190}{\ensuremath{\leftarrow}}
\providecommand{\R}{\mathbb{R}}
\providecommand{\N}{\mathbb{N}}
\providecommand{\Z}{\mathbb{Z}}
\providecommand{\Q}{\mathbb{Q}}
\providecommand{\D}{\mathbb{D}}
\providecommand{\C}{\mathbb{C}}
\providecommand{\e}{{\mathrm{e}}}
\providecommand{\vect}[1]{\overrightarrow{#1}}
\providecommand{\vectt}[1]{\overrightarrow{\mathrm{#1}}}
\providecommand{\vvt}[1]{\overrightarrow{\mathrm{#1}}}
\providecommand{\np}[2][]{{#2~#1}}
\providecommand{\vv}[1]{{\overrightarrow{#1}}}
\renewcommand{\d}{{\mathrm{d}}}
\providecommand{\degres}{{^\circ}}
\providecommand{\degre}{{^\circ}}
\providecommand{\blue}{\color{blue}}
\providecommand{\red}{\color{red}}
\providecommand{\enskip}{\quad}
\providecommand{\llbracket}{\lbrack\!\lbrack}
\providecommand{\rrbracket}{\rbrack\!\rbrack}
\providecommand{\euro}{{\text{€}}}
\providecommand{\no}{{\text{n°}}}
\providecommand{\interval}[2]{{[#1;#2]}}
\providecommand{\intervFO}[2]{{[#1;#2[}}
\providecommand{\dotfill}{\ldots}

\providecommand{\ds}{\displaystyle}
\providecommand{\textding}[1]{#1}
\providecommand{\Oij}{\(\mathrm{O};\vec{\imath},\vec{\jmath}\)}
\providecommand{\Oijk}{\(\mathrm{O};\vec{\imath},\vec{\jmath},\vec{k}\)}
\providecommand{\Ouv}{\(\mathrm{O};\vec{u},\vec{v}\)}
\providecommand{\intervFF}[2]{\ensuremath{[#1;#2]}}
\renewcommand{\euro}{\ifmmode\text{\texteuro}\else\texteuro\fi}
\DeclareUnicodeCharacter{00B0}{\ifmmode^{\circ}\else\textdegree\fi}
\DeclareUnicodeCharacter{20AC}{\ifmmode\text{\texteuro}\else\texteuro\fi}
% Énoncés : Éducation nationale. Transcription et corrigés : APMEP, Denis Vergès et contributeurs. Découpage et classement : CoopMaths, Sébastien Lozano, Éric Elter et contributeurs. Adaptation de la présentation : Maths Explorer.
% CC BY-SA : https://coopmaths.fr/www/about/
% https://coopmaths.fr/alea/static/e3c/2020/tex/e3c_2020_00_sujet59_3_cor.tex
\begin{document}
\noindent\textit{APMEP / CoopMaths. CC BY-SA. Présentation adaptée par Maths Explorer.}\par\medskip
\( f(x) = (x^2 - 2{,}5x + 1) \e^x \).

\subsection*{1.}

\paragraph{a.} \(f\) est dérivable sur \( \mathbb{R} \) comme produit de fonctions dérivables sur cet intervalle :
\begin{align*}
f'(x) &= (2x - 2{,}5) \e^x + (x^2 - 2{,}5x + 1) \e^x \\
&= \e^x (2x - 2{,}5 + x^2 - 2{,}5x + 1) \\
&= \e^x (x^2 - 0{,}5x - 1{,}5).
\end{align*}

\paragraph{b.} On sait que \( \e^x > 0 \) quel que soit le réel \( x \), donc \( f'(x) \) a le signe du trinôme \( x^2 - 0{,}5x - 1{,}5 \).

On a :
\[
\Delta = 0{,}5^2 - 4 \times (-1{,}5) = 0{,}25 + 6 = 6{,}25 = (2{,}5)^2.
\]
L'équation \( x^2 - 0{,}5x - 1{,}5 = 0 \) a donc deux racines :
\[
x_1 = \dfrac{0{,}5 + 2{,}5}{2} = 1{,}5 \quad \text{et} \quad x_2 = \dfrac{0{,}5 - 2{,}5}{2} = -1.
\]

On sait de plus que ce trinôme est positif sauf sur l'intervalle \( \,] -1 \,;\, 1{,}5[ \).

La fonction \(f\) est donc croissante sauf sur l'intervalle \( \,] -1 \,;\, 1{,}5[ \).

\subsection*{2.}

\paragraph{a.} On a :
\[
M(x \,;\, y) \in \mathcal{T} \iff y - f(0) = f'(0)(x - 0).
\]

Avec \(f(0) = 1 \times \e^0 = 1\) et \(f'(0) = -1{,}5 \times \e^0 = -1{,}5\), on a donc :
\[
M(x \,;\, y) \in \mathcal{T} \iff y - 1 = -1{,}5(x - 0) \text{ ou encore } y = -1{,}5x + 1.
\]

\paragraph{b.} Le point d'abscisse \(a\) appartenant à la fois à la courbe \(\mathcal{C}_f\) et à la tangente \(\mathcal{T}\), cette abscisse vérifie à la fois l'équation de \(f\) et celle de \(\mathcal{T}\), soit :
\[
(a^2 - 2{,}5a + 1) \e^a = -1{,}5a + 1 \text{ ou } (a^2 - 2{,}5a + 1) \e^a + 1{,}5a - 1 = 0.
\]

On peut entrer sur la calculatrice la fonction \( x \longmapsto g(x) = (x^2 - 2{,}5x + 1) \e^x + 1{,}5x - 1 \) et essayer de chercher pour quelle valeur cette fonction s'annule (en dehors de 0). On a \( g(1{,}7) \approx -0{,}42 \) et \( g(1{,}8) \approx 0{,}13 \). Donc \( 1{,}7 < a < 1{,}8 \).

\( g(1{,}77) \approx -0{,}0599 \) et \( g(1{,}78) \approx 0{,}0002 \), donc \( 1{,}77 < a < 1{,}78 \).

Donc \( a \approx 1{,}8 \) au dixième près.

\end{document}
